\SourceLeaf{266}
\Needspace{8\baselineskip}
\Article{242}
Die vorigen Sätze erfassen die Komposition der Formen in ihrer vollen Allgemeinheit. Wir gehen nun zu besonderen Anwendungen über, durch die wir ihre Reihenfolge nicht unterbrechen wollten. Zuerst greifen wir das Problem des Artikels 236 wieder auf und schränken es durch die folgenden Bedingungen ein: \textit{erstens}, die zusammenzusetzenden Formen haben dieselbe Determinante, also $d=d'$; \textit{zweitens}, $m$. $m'$ sind teilerfremd; \textit{drittens}, die gesuchte Form ist aus jeder der beiden Formen $f$, $f'$ direkt zusammengesetzt. Dann sind auch $mm$, $m'm'$ teilerfremd; somit wird der größte gemeinsame Teiler von $dm'm'$, $d'mm$, \textit{das heißt,}\ $D$, $=d=d'$, und $n=n'=1$. Die vier Größen $\mathfrak{Q}$, $\mathfrak{Q}'$, $\mathfrak{Q}''$, $\mathfrak{Q}'''$, die beliebig gewählt werden können, setzen wir jeweils gleich $=-1$, $0$, $0$, $0$. Das ist stets erlaubt, außer in dem einzigen Fall, dass $a$, $a'$, $b+b'$ zugleich $=0$ sind, den wir deshalb hier beiseitelassen. Dieser Fall kann offenbar nur bei Formen auftreten, deren Determinante ein positives Quadrat ist. Dann wird $\mu$ offenbar zum größten gemeinsamen Teiler von $a$, $a'$, $b+b'$; die Zahlen $\mathfrak{P}'$, $\mathfrak{P}''$, $\mathfrak{P}'''$ müssen so gewählt werden, dass
\[
\mathfrak{P}'a+\mathfrak{P}''a'+\mathfrak{P}'''(b+b')=\mu
\]
während $\mathfrak{P}$ selbst völlig beliebig ist. Setzt man an der angeführten Stelle für $p$, $q$, $p'$, $q'$ usw. ihre Werte ein, so ergibt sich
\[
\begin{gathered}A=\frac{aa'}{\mu\mu},\\ B=\frac{1}{\mu}(\mathfrak{P}aa'+\mathfrak{P}'ab'+\mathfrak{P}''a'b+\mathfrak{P}'''(bb'+D))\end{gathered}
\]
und $C$ kann durch die Gleichung $AC=BB-D$ bestimmt werden, sofern nicht $a$ null ist und zugleich $a'=0$ gilt.

Bei dieser Lösung hängt also der Wert von $A$ nicht von den Werten der Zahlen $\mathfrak{P}$, $\mathfrak{P}'$, $\mathfrak{P}''$, $\mathfrak{P}'''$ ab (die auf unendlich viele verschiedene Arten bestimmt werden können). Dagegen nimmt $B$ andere Werte an, wenn man diesen Zahlen andere Werte gibt, und es lohnt sich zu untersuchen, wie alle Werte von $B$ miteinander zusammenhängen. Dazu bemerken wir:

I. Wie auch $\mathfrak{P}$, $\mathfrak{P}'$, $\mathfrak{P}''$, $\mathfrak{P}'''$ bestimmt werden, die daraus entstehenden Werte von $B$ sind sämtlich modulo $A$ kongruent. Wir nehmen Folgendes an. Die Festsetzung
\[
\begin{gathered}\mathfrak{P}=\mathfrak{p},\\ \mathfrak{P}'=\mathfrak{p}',\\ \mathfrak{P}''=\mathfrak{p}'',\\ \mathfrak{P}'''=\mathfrak{p}''',\\\text{ergibt } B=\mathfrak{B}\end{gathered}
\]
und die Festsetzung
\[
\begin{gathered}\mathfrak{P}=\mathfrak{p}+\mathfrak{d},\\ \mathfrak{P}'=\mathfrak{p}'+\mathfrak{d}',\\ \mathfrak{P}''=\mathfrak{p}''+\mathfrak{d}''.\\ \mathfrak{P}'''=\mathfrak{p}'''+\mathfrak{d}''',\\\text{ergibt } B=\mathfrak{B}+\mathfrak{D}\end{gathered}
\]
\SourceLeaf{267}
Dann gilt also
\[
\begin{gathered}a\mathfrak{d}'+a'\mathfrak{d}''+(b+b')\mathfrak{d}'''=0,\\
aa'\mathfrak{d}+ab'\mathfrak{d}'+a'b\mathfrak{d}''+(bb'+D)\mathfrak{d}'''=\mu\mathfrak{D}\end{gathered}
\]
Multipliziert man die erste Seite der letzteren Gleichung mit $a\mathfrak{p}'+a'\mathfrak{p}''+(b+b')\mathfrak{p}'''$, die zweite mit $\mu$, und zieht man vom ersten Produkt die Größe
\[
(ab'\mathfrak{p}'+a'b\mathfrak{p}''+(bb'+D)\mathfrak{p}''')(a\mathfrak{d}'+a'\mathfrak{d}''+(b+b')\mathfrak{d}''')
\]
ab, die wegen der ersten Gleichung offenbar $=0$ ist, so erhält man nach dem Ausmultiplizieren und Wegheben der sich aufhebenden Glieder
\[
\begin{aligned}aa'(&\mu\mathfrak{d}\\&+((b'-b)\mathfrak{p}''+c'\mathfrak{p}''')\mathfrak{d}'\\&+((b-b')\mathfrak{p}'+c\mathfrak{p}''')\mathfrak{d}''\\&-(c'\mathfrak{p}'+c\mathfrak{p}'')\mathfrak{d}''')=\mu\mu\mathfrak{D}\end{aligned}
\]
Damit ist offenbar $\mu\mu\mathfrak{D}$ durch $aa'$, oder $\mathfrak{D}$ durch $\frac{aa'}{\mu\mu}$, \textit{das heißt,}\ durch $A$ teilbar, und
\[
\mathfrak{B}\equiv\mathfrak{B}+\mathfrak{D}\ (\mathrm{mod.}\ A)
\]

II. Wenn die Werte $\mathfrak{p}$, $\mathfrak{p}'$, $\mathfrak{p}''$ $\mathfrak{p}'''$ der Zahlen $\mathfrak{P}$, $\mathfrak{P}'$, $\mathfrak{P}''$, $\mathfrak{P}'''$ den Wert $B=\mathfrak{B}$ ergeben, so lassen sich andere Werte dieser Zahlen finden, durch die $B$ jeden vorgegebenen, zu $\mathfrak{B}$ modulo $A$ kongruenten Wert erhält, etwa $\mathfrak{B}+kA$. Zuerst bemerken wir, dass die vier Zahlen $\mu$, $c$, $c'$, $b-b'$ keinen gemeinsamen Teiler haben können. Denn hätten sie einen, so würde er die sechs Zahlen $a$, $a'$, $b+b'$, $c$, $c'$, $b-b'$ teilen, also sowohl $a$, $2b$, $c$ als auch $a'$, $2b'$, $c'$ und deshalb auch $m$, $m'$, die nach Voraussetzung teilerfremd sind. Somit lassen sich vier ganze Zahlen $h$, $h'$, $h''$, $h'''$ so angeben, dass
\[
h\mu+h'c+h''c'+h'''(b-b')=1
\]
Setzt man dann
\[
\begin{aligned}
kh&=\mathfrak{d},\\k(h''(b+b')-h'''a')&=\mu\mathfrak{d}'\\
k(h'(b+b')+h'''a)&=\mu\mathfrak{d}'',\\-k(h'a'+h''a)&=\mu\mathfrak{d}'''
\end{aligned}
\]
so sind $\mathfrak{d}$, $\mathfrak{d}'$, $\mathfrak{d}''$, $\mathfrak{d}'''$ offenbar ganzzahlig. Ferner bestätigt man leicht, dass
\[
a\mathfrak{d}'+a'\mathfrak{d}''+(b+b')\mathfrak{d}'''=0
\]
\[
\begin{aligned}aa'\mathfrak{d}+ab'\mathfrak{d}'+a'b\mathfrak{d}''+(bb'+D)\mathfrak{d}'''\\&=\frac{aa'k}{\mu}(\mu h+ch'+c'h''+(b-b')h''')\\&=\mu kA\end{aligned}
\]
Aus der ersten Gleichung folgt, dass auch $\mathfrak{p}+\mathfrak{d}$, $\mathfrak{p}'+\mathfrak{d}'$, $\mathfrak{p}''+\mathfrak{d}''$, $\mathfrak{p}'''+\mathfrak{d}'''$ Werte der Zahlen $\mathfrak{P}$, $\mathfrak{P}'$, $\mathfrak{P}''$, $\mathfrak{P}'''$ sind; aus der zweiten, dass diese Werte $B=\mathfrak{B}+kA$ ergeben. \textit{Q. E. D.} ---
\SourceLeaf{268}
Daraus ist klar, dass $B$ stets so bestimmt werden kann, dass es einschließlich der Grenzen zwischen $0$ und $A-1$ liegt, wenn $A$ positiv ist, oder zwischen $0$ und $-A-1$, wenn $A$ negativ ist.

\Needspace{8\baselineskip}
\Article{243}
Aus den Gleichungen
\[
\begin{gathered}\mathfrak{P}'a+\mathfrak{P}''a'+\mathfrak{P}'''(b+b')=\mu,\\
B=\frac{1}{\mu}(\mathfrak{P}aa'+\mathfrak{P}'ab'+\mathfrak{P}''a'b+\mathfrak{P}'''(bb'+D))\end{gathered}
\]
folgt
\[
\begin{aligned}B&=b+\frac{a}{\mu}(\mathfrak{P}a'+\mathfrak{P}'(b'-b)-\mathfrak{P}'''c)\\&=b'+\frac{a'}{\mu}(\mathfrak{P}a+\mathfrak{P}''(b-b')-\mathfrak{P}'''c')\end{aligned}
\]
also
\[
B\equiv b\ (\mathrm{mod.}\ \tfrac{a}{\mu})\quad\text{und}\quad B\equiv b'\ (\mathrm{mod.}\ \tfrac{a'}{\mu})
\]
Sind $\frac{a}{\mu}$, $\frac{a'}{\mu}$ teilerfremd, so liegt zwischen $0$ und $A-1$ (oder zwischen $0$ und $-A-1$, wenn $A$ negativ ist) nur eine einzige Zahl, die modulo $\frac{a}{\mu}$ $\equiv b$ und $\equiv b'$ modulo $\frac{a'}{\mu}$ ist. Setzt man diese Zahl $=B$ und $\frac{BB-D}{A}=C$, so ist offenbar $(A,B,C)$ aus den Formen $(a,b,c)$, $(a',b',c')$ zusammengesetzt. In diesem Fall braucht man also zur Bestimmung der zusammengesetzten Form die Zahlen $\mathfrak{P}$, $\mathfrak{P}'$, $\mathfrak{P}''$, $\mathfrak{P}'''$ nicht mehr zu berücksichtigen\HistoricalNote{*)}{Das erreicht man stets mithilfe der Kongruenzen
\[
\frac{aB}{\mu}\equiv\frac{ab'}{\mu},\quad\frac{a'B}{\mu}\equiv\frac{a'b}{\mu},\quad\frac{(b+b')B}{\mu}\equiv\frac{bb'+D}{\mu}\ (\mathrm{mod.}\ A)
\]
}. Sucht man \textit{zum Beispiel}\ eine aus den Formen $(10,3,11)$, $(15,2,7)$ zusammengesetzte Form, so sind $a$, $a'$, $b+b'$ jeweils $=10,15,5$; $\mu=5$; also $A=6$; $B\equiv3\ (\mathrm{mod.}\ 2)$ und $\equiv2\ (\mathrm{mod.}\ 3)$, woraus $B=5$ und $(6,5,21)$ als gesuchte Form folgen. --- Die Bedingung, dass $\frac{a}{\mu}$, $\frac{a'}{\mu}$ teilerfremd sind, ist genau dazu gleichwertig, dass die beiden Zahlen $a$, $a'$ keinen größeren gemeinsamen Teiler haben als die drei $a$, $a'$, $b+b'$; oder, was auf dasselbe hinausläuft, dass der größte gemeinsame Teiler von $a$, $a'$ auch die Zahl $b+b'$ teilt. Besonders seien die folgenden Fälle beachtet:

1) Es seien zwei Formen $(a,b,c)$, $(a',b',c')$ derselben Determinante $D$ gegeben, sodass der größte gemeinsame Teiler von $a$, $2b$, $c$ zum größten gemeinsamen Teiler von $a'$, $2b'$, $c'$ teilerfremd ist und $a$ zu $a'$ teilerfremd ist. Die aus ihnen zusammengesetzte Form $(A,B,C)$ findet man, indem man $A=aa'$, $B\equiv b\ (\mathrm{mod.}\ a)$ und $\equiv b'\ (\mathrm{mod.}\ a')$, $C=\frac{BB-D}{A}$ setzt. Dieser Fall liegt stets vor, wenn eine der zusammenzusetzenden Formen die Hauptform ist, etwa $a=1$, $b=0$, $c=-D$. Dann gilt $A=a'$, und $B$ kann
\SourceLeaf{269}
$=b'$ gesetzt werden, woraus $C=c'$ folgt; also \textit{ist die aus der Hauptform und irgendeiner anderen Form derselben Determinante zusammengesetzte Form diese andere Form selbst.}

2) Sollen zwei \textit{entgegengesetzte}, eigentlich primitive Formen, etwa $(a,b,c)$ und $(a,-b,c)$, zusammengesetzt werden, so ist $\mu=a$. Daraus erkennt man leicht, dass die Hauptform $(1,0,-D)$ aus ihnen zusammengesetzt ist.

3) Es seien beliebig viele eigentlich primitive Formen $(a,b,c)$, $(a',b',c')$, $(a'',b'',c'')$ usw. derselben Determinante $D$ gegeben, deren erste Koeffizienten $a$, $a'$, $a''$ usw. paarweise teilerfremd sind. Die aus ihnen allen zusammengesetzte Form $(A,B,C)$ wird gefunden. Dazu setzt man $A$ gleich dem Produkt aller $a$, $a'$, $a''$ usw.; $B$ kongruent zu $b$, $b'$, $b''$ usw. jeweils modulo $a$, $a'$, $a''$ usw.; $C=\frac{BB-D}{A}$. Denn man erkennt leicht, dass aus den beiden Formen $(a,b,c)$, $(a',b',c')$ die Form $(aa',B,\frac{BB-D}{aa'})$ zusammengesetzt ist; aus dieser und $(a'',b'',c'')$ die Form $(aa'a'',B,\frac{BB-D}{aa'a''})$ usw. --- Umgekehrt:

4) Ist eine eigentlich primitive Form $(A,B,C)$ der Determinante $D$ gegeben, zerlegt man den Koeffizienten $A$ in beliebig viele paarweise teilerfremde Faktoren $a$, $a'$, $a''$ usw., wählt Zahlen $b$, $b'$, $b''$ usw. gleich $B$ oder wenigstens dazu kongruent jeweils modulo $a$, $a'$, $a''$ usw. und setzt $ac=bb-D$, $a'c'=b'b'-D$, $a''c''=b''b''-D$ usw., so ist die Form $(A,B,C)$ aus den Formen $(a,b,c)$, $(a'\ b'\ c')$, $(a''\ b''\ c'')$ zusammengesetzt oder \textit{in diese Formen zerlegbar}. Ohne Mühe beweist man, dass derselbe Satz weiterhin gilt, selbst wenn die Form $(A,B,C)$ uneigentlich primitiv oder abgeleitet ist. Auf diese Weise lässt sich also jede Form in andere derselben Determinante zerlegen, deren erste Koeffizienten sämtlich Primzahlen oder Primzahlpotenzen sind. Eine solche Zerlegung kann oft nützlich angewendet werden, wenn aus mehreren gegebenen Formen eine zusammengesetzt werden soll. Sucht man \textit{zum Beispiel}\ die aus den Formen $(3,1,134)$, $(10,3,41)$, $(15,2,27)$ zusammengesetzte Form, so zerlege man die zweite in $(2,1,201)$, $(5,-2,81)$, die dritte in $(3,-1,134)$, $(5,2,81)$. Offenbar ist die aus den fünf Formen $(3,1,134)$, $(2,1,201)$, $(5,-2,81)$, $(3,-1,134)$, $(5,2,81)$ zusammengesetzte Form, in welcher Reihenfolge sie auch genommen werden, auch aus den drei gegebenen Formen zusammengesetzt. Durch Komposition der ersten mit der vierten entsteht aber die Hauptform $(1,0,401)$; dieselbe entsteht durch Komposition der dritten mit der fünften; somit ergibt die Komposition aller Formen die Form $(2,1,201)$.

\SourceLeaf{270}
5) Wegen des Nutzens dieser Methode lohnt es sich, sie noch weiter zu erläutern. Aus der vorigen Bemerkung ist klar, dass sich das Problem, beliebig viele gegebene eigentlich primitive Formen derselben Determinante zusammenzusetzen, auf die Komposition von Formen zurückführen lässt, deren erste Koeffizienten Primzahlpotenzen sind (denn eine Primzahl kann als die erste Potenz ihrer selbst betrachtet werden). Deshalb ist vor allem der Fall zu betrachten, dass zwei eigentlich primitive Formen $(a,b,c)$, $(a',b',c')$ zusammengesetzt werden sollen, bei denen $a$ und $a'$ Potenzen \textit{derselben} Primzahl sind. Es sei also $a=h^{\varkappa}$, $a'=h^{\lambda}$, wobei $h$ eine Primzahl bezeichnet, und wir nehmen an, was erlaubt ist, dass $\varkappa$ nicht kleiner als $\lambda$ ist. Dann ist $h^{\lambda}$ der größte gemeinsame Teiler von $a$, $a'$. Teilt dieser außerdem $b+b'$, so liegt der zu Beginn dieses Artikels betrachtete Fall vor, und $(A,B,C)$ ist aus den gegebenen Formen zusammengesetzt, wenn man $A=h^{\varkappa-\lambda}$, $B\equiv b\ (\mathrm{mod.}\ h^{\varkappa-\lambda})$ und $\equiv b'\ (\mathrm{mod.}\ 1)$ setzt, wobei die letztere Bedingung offenbar weggelassen werden kann; $C=\frac{BB-D}{A}$. Teilt dagegen $h^{\lambda}$ die Zahl $b+b'$ nicht, so muss der größte gemeinsame Teiler dieser Zahlen selbst eine Potenz von $h$ sein. Er sei also $=h^{\nu}$, dann ist $\nu<\lambda$ (man muss $\nu=0$ setzen. wenn etwa $h^{\lambda}$ und $b+b'$ teilerfremd sind). Bestimmt man also $\mathfrak{P}'$, $\mathfrak{P}''$, $\mathfrak{P}'''$ so, dass
\[
\mathfrak{P}'h^{\varkappa}+\mathfrak{P}''h^{\lambda}+\mathfrak{P}'''(b+b')=h^{\nu}
\]
wählt aber $\mathfrak{P}$ beliebig, so ist die Form $(A,B,C)$ aus den gegebenen Formen zusammengesetzt, wenn man
\[
\begin{gathered}A=h^{\varkappa+\lambda-2\nu},\\
B=b+h^{\varkappa-\nu}(\mathfrak{P}h^{\lambda}-\mathfrak{P}'(b-b')-\mathfrak{P}'''c),\\
C=\frac{BB-D}{A}\end{gathered}
\]
setzt. Man erkennt aber leicht, dass in diesem Fall auch $\mathfrak{P}'$ beliebig gewählt werden kann. Setzt man daher $\mathfrak{P}=\mathfrak{P}'=0$, so wird
\[
B=b-\mathfrak{P}'''ch^{\varkappa-\nu}
\]
oder allgemein
\[
B=kA+b-\mathfrak{P}'''ch^{\varkappa-\nu}
\]
wobei $k$ eine beliebige Zahl bezeichnet (voriger Artikel). In diese sehr einfache Formel geht allein $\mathfrak{P}'''$ ein; es ist ein Wert des Ausdrucks $\frac{h^{\nu}}{b+b'}\ (\mathrm{mod.}\ h^{\lambda})$\HistoricalNote{*)}{oder des Ausdrucks $\frac{1}{\frac{b+b'}{h^{\nu}}}\ (\mathrm{mod.}\ h^{\lambda-\nu})$, woraus $B\equiv b-\frac{ch^{\varkappa-\nu}}{\frac{b+b'}{h^{\nu}}}\equiv\frac{(D+bb'):h^{\nu}}{(b+b'):h^{\nu}}\ (\mathrm{mod.}\ A)$ folgt}. Sucht man \textit{zum Beispiel}\ die aus $(16,3,19)$ und $(8,1,37)$ zusammengesetzte Form, so ist $h=2$, $\varkappa=4$, $\lambda=3$, $\nu=2$.
\SourceLeaf{271}
Also ist $A=8$, und $\mathfrak{P}'''$ ein Wert des Ausdrucks $\frac{4}{4}\ (\mathrm{mod.}\ 8)$, etwa $1$, woraus $B=8k-73$ folgt. Setzt man also $k=9$, so ist $B=-1$ und $C=37$, oder $(8,-1,37)$ die gesuchte Form.

Sind also beliebig viele Formen gegeben, deren erste Koeffizienten sämtlich Primzahlpotenzen sind, so muss man nachsehen, ob die ersten Koeffizienten einiger Formen Potenzen \textit{derselben} Primzahl sind, und diese jeweils nach der eben gegebenen Regel miteinander zusammensetzen. So entstehen Formen, deren erste Koeffizienten weiterhin Primzahlpotenzen sind, aber Potenzen durchweg verschiedener Primzahlen; die aus ihnen zusammengesetzte Form kann daher nach der dritten Bemerkung bestimmt werden. \textit{Zum Beispiel:}\ Sind die Formen $(3,1,47)$, $(4,0,35)$, $(5,0,28)$, $(16,2,9)$, $(9,7,21)$, $(16,6,11)$ gegeben, so entsteht aus der ersten und fünften die Form $(27,7,7)$; aus der zweiten und vierten entsteht $(16,-6,11)$, aus dieser und der sechsten $(1,0,140)$, die weggelassen werden kann. Es bleiben also $(5,0,28)$, $(27,7,7)$, aus denen $(135,-20,4)$ entsteht; stattdessen kann man die eigentlich äquivalente Form $(4,0,35)$ nehmen. Diese ist also das Ergebnis der Komposition der sechs gegebenen Formen.

Aus dieser Quelle lassen sich noch andere in der Anwendung nützliche Verfahren schöpfen. Doch um nicht zu weit auszuholen, lassen wir eine ausführlichere Behandlung weg und eilen zu anderen, schwierigeren Dingen.
